\chapter{Introdução}\label{cap:introducao}

\TextoModelo{Apresente o tema, delimite o problema geológico, sintetize o estado do conhecimento, justifique o estudo e antecipe a abordagem utilizada. A introdução deve conduzir naturalmente aos objetivos.}

Como exemplo de citação no sistema autor-data, uma referência pode ser integrada ao texto por meio de \citeonline{hasui2012} ou apresentada entre parênteses \cite{press2006}. Substitua estas referências pelas fontes efetivamente utilizadas.

\section{Objetivos}\label{sec:objetivos}

\subsection{Objetivo geral}

\TextoModelo{Declare o resultado científico principal que o trabalho pretende alcançar.}

\subsection{Objetivos específicos}

\begin{itemize}
  \item \TextoModelo{caracterizar as unidades geológicas da área de estudo;}
  \item \TextoModelo{descrever os procedimentos de campo, laboratório e gabinete;}
  \item \TextoModelo{integrar e interpretar os dados obtidos;}
  \item \TextoModelo{discutir as implicações geológicas dos resultados.}
\end{itemize}

\section{Justificativa}

Apesar de serem menos estudados quando comparado aos deslizamentos terrestres, os movimentos de massa em ambiente marinho também são importantes em termos sociais, econômicos e ecológicos (Clare et al., 2018), uma vez que são fontes potenciais a gerar danos (Vanneste et al., 2014). Frequentemente, esses deslizamentos submarinos geram fluxos de dezenas a centenas de quilômetros, capazes de danificar a infraestrutura estratégica no fundo do mar, como cabos de telecomunicações, plataformas de produção e dutos de hidrocarbonetos (Pope et al., 2016), os quais estão presentes na região norte da Bacia de Santos.

Sob soterramento, os depósitos de deslizamentos subaquáticos são elementos importantes nos sistemas petrolíferos, pois influenciam a distribuição dos reservatórios (Kneller et al., 2016), podem atuar como selos (Cardona et al., 2016) e também como potenciais reservatórios (Meckel, 2011).

Um dos maiores desafios de estudos sobre transportes de massa submarinos é a sua natureza, em mar aberto e longe da costa, o que dificulta a aquisição dos dados bem como a falta de informação consistente, por vezes, de baixa resolução (Clare et al., 2018). Assim, convergir conhecimentos da atividade sismoestratigráfica e sismológica, visando a identificação e caracterização de movimentos de massa, dará suporte para avançar no entendimento de processos atuantes na margem continental sudeste e, consequentemente, no planejamento de ações da indústria do petróleo na implantação e monitoramento de infraestrutura de exploração e produção de hidrocarbonetos.
